\documentclass[12pt,a4paper]{ctexart} % 使用 ctexart，直接支持中文
\usepackage[margin=2.2cm]{geometry} % 设置页边距
\usepackage{xcolor} % 颜色支持
\usepackage{hyperref} % 超链接
\usepackage{subfiles} % 允许课程文件既能被主文件引入，也能单独编译
\usepackage{enumitem} % 更紧凑的列表
\usepackage{longtable} % 跨页表格
\usepackage{array} % 表格列格式增强
\usepackage{fvextra} % 更稳定地展示可自动换行的代码块

\hypersetup{
  colorlinks=true,
  linkcolor=blue!60!black,
  urlcolor=blue!60!black
}

\setlist{nosep} % 减少列表上下留白
\setlength{\parindent}{2em} % 首行缩进
\setlength{\parskip}{0.4em} % 段间距

% 统一定义代码块环境，课程文件里直接复用
\DefineVerbatimEnvironment{CodeBlock}{Verbatim}{
  fontsize=\small,
  frame=single,
  breaklines=true
}
\DefineShortVerb{\`}

\title{LaTeX / Overleaf 零基础动手课程}
\author{latexlearningCODEX}
\date{\today}

\begin{document}

\maketitle

\section*{这份项目怎么用}

这不是一本偏理论的 LaTeX 教材，而是一套“边改边编译”的练习型项目。
建议你先编译本文件，看清课程顺序；然后打开对应的 `examples/` 示例文件，一边修改，一边重新编译观察结果。

\section*{学习节奏建议}

\begin{enumerate}
  \item 先原样编译一次，不要急着改。
  \item 每次只改一个点，例如标题、一个命令、一个公式或一个表格单元格。
  \item 修改后立刻重新编译。
  \item 把注意力放在“哪一行代码变化，导致 PDF 哪一处变化”。
\end{enumerate}

\section*{文件导航}

\begin{longtable}{>{\raggedright\arraybackslash}p{0.17\linewidth} >{\raggedright\arraybackslash}p{0.33\linewidth} >{\raggedright\arraybackslash}p{0.40\linewidth}}
\textbf{课程} & \textbf{讲义文件} & \textbf{建议修改的示例文件} \\
\hline
第 1 课 & lessons/lesson01\_first\_document.tex & examples/lesson01\_hello.tex \\
第 2 课 & lessons/lesson02\_document\_structure.tex & examples/lesson02\_structure.tex \\
第 3 课 & lessons/lesson03\_sections\_and\_text.tex & examples/lesson03\_sections\_text.tex \\
第 4 课 & lessons/lesson04\_math\_basics.tex & examples/lesson04\_math.tex \\
第 5 课 & lessons/lesson05\_figures.tex & examples/lesson05\_figures.tex \\
第 6 课 & lessons/lesson06\_tables.tex & examples/lesson06\_tables.tex \\
第 7 课 & lessons/lesson07\_references\_and\_bib.tex & examples/lesson07\_refs\_bib.tex \\
第 8 课 & lessons/lesson08\_final\_mini\_report.tex & examples/final\_report\_template.tex \\
\end{longtable}

\tableofcontents
\newpage

\subfile{lessons/lesson01_first_document}
\subfile{lessons/lesson02_document_structure}
\subfile{lessons/lesson03_sections_and_text}
\subfile{lessons/lesson04_math_basics}
\subfile{lessons/lesson05_figures}
\subfile{lessons/lesson06_tables}
\subfile{lessons/lesson07_references_and_bib}
\subfile{lessons/lesson08_final_mini_report}

\end{document}
